\section{Encoding complete choiceless computations}
\label{sec:coding}

The preceding transfer theorem concerns hereditarily finite sets with small
supports. We now bring an arbitrary fixed choiceless computation into that
setting. There are two separate points to prove: all constituents of its
encoding lie in a polynomial-size invariant family, and its states can be
defined by counting formulas whose number of variable names stays bounded
throughout the computation. For our two inputs, which have the same size,
this will give one common $C^m$ sentence for each stage's halt test and output test, exact
on their supported domains through the respective halt stages, with $m$
depending only on the program.

\subsection{The interpretation characterization}

We use the binary-vocabulary formulation of the interpretation characterization
of choiceless polynomial time with counting, due to Grohe, Schweitzer and
Wiebking~\cite[Section~6 and Theorem~18, arXiv version~1]{gsw2020}, building on
Gr\"adel, Pakusa, Schalth\"ofer and Kaiser~\cite[Theorem~1]{gkps2015}.
Here counting includes the cardinality operation on hereditarily finite sets.
We spell out the exact direction and interpretation convention used below.

An interpretation from one fixed binary relational vocabulary to another
consists of fixed, parameter-free formulas
\[
 \delta(x_1,x_2),\qquad
 \eta(x_1,x_2,y_1,y_2),\qquad
 \rho_R(x_1,x_2,y_1,y_2)
\]
in first-order logic with the H\"artig quantifier. This quantifier compares
the cardinalities of two sets of elements defined by formulas; the two
formulas may have free variables in addition to their respective counted
variables. On a finite structure with universe $V$, the interpretation first
forms
\[
 D=\{(a_1,a_2)\in V^2:\delta(a_1,a_2)\}.
\]
Let $\sim$ be the least equivalence relation on $D$ containing every pair
$(\bar a,\bar b)$ satisfying $\eta(\bar a,\bar b)$. The new universe is
$D/{\sim}$, and
\[
 R([\bar a],[\bar b])
 \quad\Longleftrightarrow\quad
 \text{some }\bar a'\sim\bar a,\ \bar b'\sim\bar b
 \text{ satisfy }\rho_R(\bar a',\bar b').
\]
Thus the interpretation is defined even when $\eta$ itself is not an
equivalence relation or $\rho_R$ is not constant on classes. This is the
convention in~\cite{gsw2020}.

An interpretation program has one such initialization interpretation, one
step interpretation, and fixed halt and output sentences in first-order
logic with the H\"artig quantifier. Write $A_0=A$ for the input,
$A_1=I_{\rm init}(A_0)$, and $A_{j+1}=I_{\rm step}(A_j)$ for $j\ge1$.
The halt sentence is first tested at $A_1$. A deciding program halts at a
stage $h\ge1$, outputs the truth value of its output sentence there, and has
\[
 h+\sum_{j=1}^{h}|V(A_j)|
\]
bounded by a fixed polynomial in $|V(A)|$. The characterization says, in the
direction needed here, that every property of finite binary relational
structures definable in full choiceless polynomial time with counting is
decided by such a program. The absence of parameters means that the fixed
interpretations have no additional input-dependent distinguished tuple;
it does not prohibit free tuple variables in their defining formulas.

Fix a deciding interpretation program $P$. On our two inputs of the same
size $N$, choose a common integer $B\ge2$ bounding $N$, both halt stages,
and every state size up to halt. We may take $B$ to be a fixed polynomial
bound depending only on $P$.

\subsection{A polynomial hereditary encoding}

For objects $a,b$ in the universe of hereditarily finite sets over the input
atoms, use the Kuratowski pair
\[
 \langle a,b\rangle=\{\{a\},\{a,b\}\}.
\]
This coding is injective in the ordered pair, including when $a$ or $b$ is
an atom. Indeed the intersection of its members is $\{a\}$, while their
union is $\{a,b\}$; after recovering $a$, the latter set determines $b$,
also when $a=b$.

Encode the vertices of $A_0$ by the input atoms themselves. Suppose the
vertices of $A_{j-1}$ have already been encoded injectively. A vertex $C$
of $A_j$ is an equivalence class of domain pairs from $A_{j-1}$. Define
its payload and code by
\[
 S_C=\{\langle a_1,a_2\rangle:
       (a_1,a_2)\text{ is a pair of previous codes representing }C\},
 \qquad
 e_j(C)=\langle [j],S_C\rangle,
\]
where $[j]=\{[0],\ldots,[j-1]\}$ is the von Neumann ordinal. Distinct
classes have disjoint, nonempty sets of representing pairs, so their
payloads, and hence their codes, are distinct. The payload records the
whole class, so its construction does not choose a representative. The ordinal tag distinguishes vertices created at
different stages, and stage-zero codes are atoms rather than sets.

\begin{lemma}\label{lem:trace-size}
For either input, the input atoms, all vertex codes up to halt, and all
objects recursively occurring as members of those codes belong to an
$H$-invariant transitive family of at most $9B^3$ objects.
\end{lemma}
\begin{proof}
At a single stage there are at most $B^2$ domain pairs. For each pair,
its Kuratowski code and its two inner sets add at most three objects.
There are at most $B$ classes; their payloads add at most $B$ objects,
and the tagged vertex pairs add at most $3B$ further objects. Previous
vertex codes and their constituents have already been counted. Across
all stages, the ordinal tags and their constituents add at most $B+1$
objects. Including the at most $B$ input atoms, the total is at most
\[
 B(3B^2+4B)+2B+1\le 9B^3
\]
for $B\ge2$. Taking the recursive membership closure therefore gives
the stated transitive family.

Every input automorphism preserves truth of the fixed interpretation
formulas, including the H\"artig quantifier. It therefore permutes the
domain pairs, their generated equivalence classes, and all output
relations. Inductively it permutes the vertex codes: it acts on a pair
or set by acting on its constituents, and fixes every pure ordinal.
It consequently preserves the family and its membership closure.
In particular the input action of $H$ preserves this family.
\end{proof}

Since $B$ is polynomial in $N$, fix an integer $q\ge1$, depending only on
$P$, such that $9B^3\le N^q$ for all sufficiently large $N$. Every
$H$-orbit in the transitive trace family has at most $N^q$ objects.
Corollary~\ref{cor:transitive-support} therefore gives $K$-supports of the
common length $s$ for every code and every recursive constituent. Thus
the entire family lies in the corresponding supported domain
$\mathcal D_b$, including the payloads, all pair wrappers, and the ordinal
tags. No bound on the rank of a code has been imposed: its nesting depth
may grow throughout the computation.

\subsection{Exact formulas for the encoded states}

All formulas in this subsection are evaluated in $\mathcal D_b$ with its
membership relation and its atomic input relations. Its domain is
transitive, contains all input atoms and all pure hereditarily finite sets,
and consists of atoms and sets. Consequently
\[
 \mathsf{Atom}(x):=\mathsf{Ed}(x,x)\vee\mathsf{Cf}(x,x),
 \qquad \mathsf{Set}(x):=\neg\mathsf{Atom}(x)
\]
distinguish these two kinds of objects.

Let $U_0(x)=\mathsf{Atom}(x)$ and let the stage-zero relation formulas be
the input relations. Suppose that the formulas for the preceding stage
are exact: $U(x)$ selects precisely its vertex codes, and its relation
formulas describe exactly the state on those codes. Translate a fixed
first-order formula with H\"artig quantifiers by replacing state relation
atoms by their defining formulas, retaining equality, and restricting
ordinary quantifiers to $U$. In particular, equality is equality of the
injective vertex codes, including for repeated variable arguments.
Replace a H\"artig comparison of $\chi_1(x)$ and $\chi_2(y)$ by
\begin{equation}\label{eq:hartig-translation}
 \bigvee_{n=0}^{B}
 \left[
  \bigl(\exists^{=n}x\,(U(x)\wedge\chi_1^*)\bigr)
  \wedge
  \bigl(\exists^{=n}y\,(U(y)\wedge\chi_2^*)\bigr)
 \right].
\end{equation}
Both counted sets have size at most $B$, so induction on formulas proves
this translation exact, also when the state is empty. The superscript
$*$ will denote this translation.

For the initialization or step interpretation appropriate to the new
stage, put
\[
 \mathsf{Dom}(\bar a):=U(a_1)\wedge U(a_2)\wedge\delta^*(\bar a),
 \qquad \bar a=(a_1,a_2).
\]
On two domain pairs define
\[
 \mathsf{Link}(\bar a,\bar b):=
 \mathsf{Dom}(\bar a)\wedge\mathsf{Dom}(\bar b)\wedge
 \bigl(\bar a=\bar b\vee\eta^*(\bar a,\bar b)
                         \vee\eta^*(\bar b,\bar a)\bigr),
\]
where tuple equality abbreviates the two coordinate equalities. Define
finite formulas recursively by
\[
 \begin{split}
 \mathsf{Reach}_0(\bar a,\bar b)&:=\mathsf{Link}(\bar a,\bar b),\\
 \mathsf{Reach}_{t+1}(\bar a,\bar b)&:=
 \mathsf{Reach}_t(\bar a,\bar b)\vee
 \exists\bar d\,
  (\mathsf{Reach}_t(\bar a,\bar d)\wedge
                         \mathsf{Link}(\bar d,\bar b)).
 \end{split}
\]
This expresses paths of length at most $t+1$ in the reflexive symmetric
link graph. There are at most $B^2$ domain pairs, so
$\mathsf{Reach}_{B^2}$ is exactly the generated equivalence relation,
and is false if either endpoint is outside the domain.

The generated quotient is now defined. We next recognize its hereditary
codes; transitivity will ensure that membership tests identify the actual
sets, not merely their visible members. Write
\[
 \mathsf{Sing}(u,a):=\mathsf{Set}(u)\wedge
                  \forall w\,(w\in u\leftrightarrow w=a),
\]
and
\[
 \mathsf{Double}(v,a,b):=\mathsf{Set}(v)\wedge
               \forall w\,(w\in v\leftrightarrow(w=a\vee w=b)).
\]
Then
\[
 \mathsf{OP}(z,a,b):=
 \exists u\,\exists v\,
 (\mathsf{Sing}(u,a)\wedge\mathsf{Double}(v,a,b)
                         \wedge\mathsf{Double}(z,u,v))
\]
expresses $z=\langle a,b\rangle$ whenever that pair and its wrappers
are present. Conversely every satisfying assignment gives the actual
Kuratowski pair: transitivity ensures that the universally quantified
membership tests cannot overlook a member outside the domain.
Define the ordinal formulas by
\[
 \mathsf{Ord}_i(o):=\mathsf{Set}(o)\wedge
 \forall w\left(w\in o\leftrightarrow
                       \bigvee_{0\le h<i}\mathsf{Ord}_h(w)\right).
\]
The empty disjunction for $i=0$ is false. All pure ordinals belong to
$\mathcal D_b$, and induction using transitivity shows that this formula
selects precisely $[i]$.

For stage $j\ge1$, define
\begin{align*}
 \mathsf{Payload}_j(S,\bar a):={}&
 \mathsf{Set}(S)\wedge\mathsf{Dom}(\bar a)\\[-2pt]
 &{}\wedge\forall z\left(z\in S\leftrightarrow
  \exists\bar b\,
   (\mathsf{Dom}(\bar b)\wedge
    \mathsf{Reach}_{B^2}(\bar a,\bar b)
                      \wedge\mathsf{OP}(z,b_1,b_2))\right),\\
 \mathsf{Rep}_j(x,\bar a):={}&
 \exists o\,\exists S\,
  (\mathsf{Ord}_j(o)\wedge\mathsf{Payload}_j(S,\bar a)
                               \wedge\mathsf{OP}(x,o,S)),\\
 U_j(x):={}&\exists\bar a\,\mathsf{Rep}_j(x,\bar a),\\
 R_j(x,y):={}&\exists\bar a\,\exists\bar b\,
 (\mathsf{Rep}_j(x,\bar a)\wedge\mathsf{Rep}_j(y,\bar b)
                                      \wedge\rho_R^*(\bar a,\bar b)).
\end{align*}
The stage subscripts on the formulas and on their auxiliary definitions
are syntactic indices, not additional free variables.

\begin{lemma}\label{lem:state-exactness}
At every stage through halt, $U_j$ selects exactly the vertex codes of
$A_j$, and each $R_j$ defines exactly its corresponding binary state
relation. The translated halt and output sentences are exact at every
stage $j\ge1$ through halt.
\end{lemma}
\begin{proof}
The assertion at stage zero follows from the definition of the atomic
input relations. Assume it for the preceding stage. Formula induction,
including~\eqref{eq:hartig-translation}, makes $\mathsf{Dom}$ and the
interpretation formulas exact. The link graph and its bounded path
formula therefore give exactly the required classes of domain pairs.

For a domain pair $\bar a$ representing a class $C$, every pair code
belonging to $S_C$ and all its wrappers are in $\mathcal D_b$, by
Lemma~\ref{lem:trace-size} and the support theorem. Hence the right side
of the payload biconditional selects exactly the elements of $S_C$.
The true payload exists in the domain, while any satisfying payload
has exactly those members, since the domain is transitive. Extensionality
for sets gives $S=S_C$. The ordinal and ordered-pair formulas then force
$x=e_j(C)$, with all required witnesses present. This proves the exactness
of $\mathsf{Rep}_j$ and $U_j$. If there are no domain pairs, both formulas
are false for every candidate vertex, as required.

Finally, $R_j$ holds precisely when some representatives of the two
coded classes satisfy the defining relation formula. This is the chosen
quotient semantics, so all relations are exact. Applying the same
formula translation to the fixed halt and output sentences proves the
last assertion.
\end{proof}

\subsection{A variable bound independent of the number of stages}

The formulas just constructed may be extremely long. What matters for
the transfer theorem is a different syntactic bound. For a formula
$\theta$, define
\[
 \operatorname{fw}(\theta)
   :=\max_{\xi\text{ a subformula of }\theta}|\operatorname{FV}(\xi)|.
\]
Here free variables are counted relative to the subformula $\xi$ itself.
Thus a variable bound outside $\xi$ is counted if it occurs freely inside
$\xi$.

\begin{lemma}\label{lem:substitution-width}
Suppose predicates in a relational formula $T$ are replaced by formulas
whose free variables are among their displayed argument variables and
whose free-variable measure is at most $w$. Capture-avoiding substitution
produces a formula of free-variable measure at most
$\max\{\operatorname{fw}(T),w\}$. This allows repeated arguments and
unused displayed argument variables. Moreover a formula of measure at
most $w$ can be renamed to use at most $w+1$ variable names.
\end{lemma}
\begin{proof}
Before each insertion, rename its bound variables to avoid the argument
variables and all possible capture. A subformula inside the inserted
copy has the same free variables as its original counterpart, except
that displayed argument variables are replaced by their actual variable
arguments. This replacement can identify variables but cannot increase
their number. A subformula inherited from the template gains no new
free variables: the inserted formula has free variables only among the
arguments of the atom it replaces. The claimed maximum bound follows.
The argument also covers equality atoms, which are retained as equality,
and atomic occurrences with repeated arguments.

For the renaming claim, fix a palette of $w+1$ names. Recursively process
a formula node with an injective assignment of palette names to that
node's free variables. At a Boolean connective, restrict this assignment
to each child's free variables. At a binder $Qx\,\xi$, keep the assigned
names of the free variables of the parent and choose a name for $x$
outside those names. There is a spare name because there are at most
$w$ parent free variables. If $x$ is free in $\xi$, extend the assignment
by this name; otherwise the binder is vacuous and the child keeps the
parent assignment. Recursively rename the child. Its free-variable
assignment is injective, and its size is at most $w$. Names may be reused
in subtrees where the corresponding variable is not free. This preserves
all bindings and truth, and applies equally to ordinary and exact-count
quantifiers.
\end{proof}

Apply the lemma first to the fixed formulas of $P$. Replacing a
H\"artig quantifier by~\eqref{eq:hartig-translation} may repeat a
subformula $B+1$ times but does not increase its free-variable sets with
$B$; the finitely many program formulas therefore give a uniform bound
for the translation templates. Pair coding, payloads, representatives,
and domain and relation definitions also have fixed templates, with all
free variables among their displayed arguments.

For example, the template for $\mathsf{Reach}_{t+1}(\bar a,\bar b)$
uses only the endpoint pairs $\bar a,\bar b$ and the quantified pair
$\bar d$. The inserted copy $\mathsf{Reach}_t(\bar a,\bar d)$ exposes
only its four endpoint arguments; $\bar b$ occurs in the other conjunct,
$\mathsf{Link}(\bar d,\bar b)$, and is not a free parameter of that
copy. Thus expansion takes the maximum of the previous width and the
fixed template width, rather than adding the widths at successive
levels. The same argument applies when a stage formula is inserted into
the next stage's template. The ordinal recursion has the simpler unary
interface. Taking the maximum over these finitely many template widths
and the bounds for the fixed formulas of $P$ gives a constant $w=w(P)$
for every fully expanded formula above. Unused or identified arguments
can only reduce the number of free variables.

Set $m=\max\{2,w+1\}$. Lemma~\ref{lem:substitution-width} puts all the
expanded state, halt, and output formulas in $C^m$, independently of
$N$, $B$, and the stage number. The numerical constants in the counting
quantifiers and the lengths of these formulas may depend on $B$ and the
stage. These are formulas used to compare two runs; they are not an
additional nonuniform program.
